
(3 points + 3 points)
\begin{teilaufgaben}
\item
The function $u(x,t)$ is defined on the strip
$\Omega = [0,1] \times [0,\infty)$.

The function $u(x,t)$ satisfies in $\Omega$ the equation
\[
u_{t}(x,t) = u_{xx}(x,t)
\]
and  $u(0,t) = u(1,t) = 1$ for $t \in [0,\infty)$.

One also has the initial values $u(1/3,0) = u(2/3,0) = 2$ .

Determine approximate values for $u(1/3,1)$ and $u(2/3,1)$.

Use the implicit method of Richardson with $\Delta x = 1/3$ and
$\Delta t = 1/3$.

\item
The boundary of $\Omega$ is a closed polygon with the vertices:
\[
(0,0), (8,0), (8,2), (4,2), (0,6). 
\]

The function $u(x,y)$ is defined on $\Omega$. It satisfies in $\Omega$ 
\[
\Delta u(x,y) = 1.
\]
On the boundary of $\Omega$ one has $u(x,y) = 0$. 

Determine approximate values for $u(1,1)$ and $u(7,1)$. 

Use finite volumes with two Voronoi cells.  
\end{teilaufgaben}

\begin{loesung}
\begin{teilaufgaben}
\item
Set $U_k^{j} = \tilde u(k/3,j/3)$. We apply the implicit method of Richardson with $r = 3$ and obtain. 
\begin{equation}
\begin{pmatrix*}[r]
-3 & 7 & -3 & 0 \\
0 & -3 & 7 & -3
\end{pmatrix*}
\cdot
\begin{pmatrix}
1 \\
U_1^{j+1} \\
U_2^{j+1} \\
1
\end{pmatrix}
=
\begin{pmatrix}
U_1^{j} \\
U_2^{j}
\end{pmatrix}.
\tag{\textbf{1P}}
\end{equation}
We use the initial values to obtain:
\[
\begin{pmatrix*}[r]  7 & -3 &  \\  -3 & 7 \end{pmatrix*}
\cdot
\begin{pmatrix} U_1^{1} \\  U_2^{1}  \end{pmatrix}
-
\begin{pmatrix} 3\\  3  \end{pmatrix}
=
\begin{pmatrix} 2 \\  2  \end{pmatrix}
\]
Therefore we have for the first time-step:
\begin{equation}
\tilde u(1/3,1/3) = \tilde u(2/3, 1/3) = 5/4.
\tag{\textbf{1P}}
\end{equation}
Next we have:
\[
\begin{pmatrix*}[r]  7 & -3 &  \\  -3 & 7 \end{pmatrix*}
\cdot
\begin{pmatrix} U_1^{2} \\  U_2^{2}  \end{pmatrix}
-
\begin{pmatrix} 3\\  3  \end{pmatrix}
=
\begin{pmatrix} 5/4 \\  5/4  \end{pmatrix}
\]
Therefore we have for the second time-step:
\[
\tilde u(1/3,2/3) = \tilde u(2/3, 2/3) = 17/16.
\]
Finally we have:
\[
\begin{pmatrix*}[r]  7 & -3 &  \\  -3 & 7 \end{pmatrix*}
\cdot
\begin{pmatrix} U_1^{3} \\  U_2^{3}  \end{pmatrix}
-
\begin{pmatrix} 3\\  3  \end{pmatrix}
=
\begin{pmatrix} 17/16 \\  17/16 \end{pmatrix}
\]
Therefore we have for the last time-step:
\begin{equation}
\tilde u(1/3,1) = \tilde u(2/3, 1) = 65/64.
\tag{\textbf{1P}}
\end{equation}
\item
By the finite volume method we have
\[
2/6 \cdot (\tilde u(7,1) - \tilde u(1,1))
+
\sqrt{32}/\sqrt{8} \cdot (0 - \tilde u(1,1))
+
6/1 \cdot (0 - \tilde u(1,1))
+
4/1 \cdot (0 - \tilde u(1,1))
=
16
\]
and
\[
2/1 \cdot (\tilde 0 - \tilde u(7,1))
+
4/1 \cdot (0  - \tilde u(7,1))
+
2/6 \cdot (\tilde u(1,1)- \tilde u(7,1))
+
4/1 \cdot (0 - \tilde u(7,1)) = 8,
\]
This leads to a system of two linear equations:
\begin{equation}
1/3
\cdot
\begin{pmatrix*}[r]
-37 & 2 \\
 2 & -31
\end{pmatrix*}
\cdot
\begin{pmatrix}
\tilde u(1,1)  \\
\tilde u(7,1)
\end{pmatrix}
=
\begin{pmatrix*}[r]
16 \\
 8
\end{pmatrix*}.
\tag{\textbf{2P}}
\end{equation}
Hence $\tilde u(1,1) = -252/191 \approx -1.3194$ and
$\tilde u(7,1) = -156/191 \approx -0.8168$ (\textbf{1P}).
\qedhere
\end{teilaufgaben}
\end{loesung}
